% Auto-generated body — do not edit by hand.

{\FontStep{000} Reality is not a collection of things.}

\medskip

{\FontStep{000} At least, that is how the world first appears. We point to a chair, a tree, a building, a planet, and we imagine that reality is fundamentally composed of objects. The boundaries seem obvious. The categories seem natural. The nouns appear primary.}

\medskip

{\FontStep{000} Yet every object is caught in a process. The chair is slowly becoming dust. The tree is transforming sunlight into structure. The building expands and contracts with temperature. The planet exchanges matter and energy with its surroundings. The object is never still. It is a temporary stabilization within a larger field of change.}

\medskip

{\FontStep{001} Suppose we begin not with things but with transformations.}

\medskip

{\FontStep{001} A river remains identifiable despite the continual replacement of its water. A language remains recognizable despite centuries of drift. A city persists despite the destruction and reconstruction of its buildings. Identity appears less like a static substance and more like a trajectory through a space of possible states.}

\medskip

{\FontStep{001} This shift changes how we think about memory.}

\medskip

{\FontStep{001} When people speak of preserving information, they often imagine preserving a specific object: a book, a file, a database, a recording. Yet preservation of the object does not guarantee preservation of access. A perfectly preserved document written in an unknown language is inaccessible. A perfectly preserved file stored on unreadable hardware is inaccessible. A perfectly preserved archive buried where nobody can reach it is inaccessible.}

\medskip

{\FontStep{001} Reachability matters.}

\medskip

{\FontStep{001} A memory that cannot be recovered may exist physically while being functionally absent. A path that cannot be traversed may exist geometrically while being operationally irrelevant. Preservation and accessibility are related but distinct concepts.}

\medskip

{\FontStep{002} Consider a simple example.}

\medskip

{\FontStep{002} A person struggles to recall a name. The memory has not necessarily vanished. The information may still exist somewhere within the nervous system. What is missing is not the content itself but a path leading to that content. A clue, a photograph, a familiar voice, or a fragment of conversation may suddenly restore access. The memory returns not because it was recreated but because a route became available.}

\medskip

{\FontStep{002} The same principle appears in libraries, archives, ecosystems, and computation.}

\medskip

{\FontStep{002} Large systems survive by maintaining routes between states. A repair manual is useful because it provides a path from failure to function. A road network is useful because it provides a path from one location to another. A language is useful because it provides a path from one mind to another.}

\medskip

{\FontStep{002} Every successful system is, in some sense, a machine for preserving reachable futures.}

\medskip

{\FontStep{003} Numbers:
0123456789}

\medskip

{\FontStep{003} Punctuation:
! @ \# \$ \% \textasciicircum{} \& * ( ) [ ] \{ \} < > ? / \textbackslash{} | + = - \_ : ; , .}

\medskip

{\FontStep{003} Uppercase:
ABCDEFGHIJKLMNOPQRSTUVWXYZ}

\medskip

{\FontStep{003} Lowercase:
abcdefghijklmnopqrstuvwxyz}

\medskip

{\FontStep{003} The quick brown fox jumps over the lazy dog.}

\medskip

{\FontStep{004} Pack my box with five dozen liquor jugs.}

\medskip

{\FontStep{004} Sphinx of black quartz, judge my vow.}

\medskip

{\FontStep{004} How razorback-jumping frogs can level six piqued gymnasts!}

\medskip

{\FontStep{004} Now imagine a document that slowly changes as it is read.}

\medskip

{\FontStep{004} At first the alterations are subtle. A curve becomes slightly irregular. A stem becomes slightly narrower. The typography acquires a faint instability. The reader notices nothing.}

\medskip

{\FontStep{004} Further into the text the changes accumulate. Counters shrink. Spacing drifts. The distinction between similar letters becomes less reliable. Reading requires additional effort.}

\medskip

{\FontStep{005} Still later entire glyph classes begin to mutate. Familiar symbols acquire unexpected features. Shapes from other writing systems appear. Ligatures emerge and disappear. Some characters become ambiguous.}

\medskip

{\FontStep{005} Eventually the document enters a different representational regime. The symbols remain structured, but the original alphabet can no longer be assumed. Interpretation becomes an act of reconstruction.}

\medskip

{\FontStep{005} Near the end, only fragments remain.}

\medskip

{\FontStep{005} Words become patterns.}

\medskip

{\FontStep{005} Patterns become traces.}

\medskip

{\FontStep{006} Traces become constraints.}

\medskip

{\FontStep{006} The reader is left not with a collection of preserved symbols but with the residual geometry of what once could be reached.}

\medskip

{\FontStep{006} And perhaps that is the deeper lesson.}

\medskip

{\FontStep{006} Every archive decays.}

\medskip

{\FontStep{006} Every language drifts.}

\medskip

{\FontStep{006} Every civilization transforms.}

\medskip

{\FontStep{007} Every memory fades.}

\medskip

{\FontStep{007} What survives is not necessarily the object itself.}

\medskip

{\FontStep{007} What survives is the possibility of finding a path back.}

\medskip
